% English original: sections/01-introduction.tex
\section{서론}

배틀로얄 match의 최종 순위는 하나의 정적인 결과가 아니라, 축소되는 전장
안에서 phase마다 변화하는 생존 위험이 누적된 결과이다. 팀의 탈락 위험은
팀이 어디에 서 있는지, 어떤 상대가 근처에 있는지, 수축하는 구역이 얼마나
강하게 압박하는지에 따라 오르내린다. 총 kill 수, 총 damage, 최종 이동 거리
같은 경기 후 집계는 이 과정을 시간 순서가 제거된 사후 통계로 평탄화한다.
같은 양의 damage라도 초반의 소규모 교전과 좁은 구역 안의 후반 전투에서는
다른 위험을 뜻하고, 같은 이동 거리라도 초반에 선택한 이동과 후반에 강요된
이동은 다른 노출을 뜻한다. 따라서 match 수준의 label로는 팀이 언제
위험해졌는지, 어떤 공간 조건이 원인이었는지 말할 수 없다.

본 연구는 이 과정을 \textit{PUBG: Battlegrounds}에서 다룬다. PUBG는 이
질문에 구조적으로 적합하다. 한 match는 최대 100명의 플레이어를 넓은 map에
배치하고, 번호가 매겨진 \textit{phase}를 거치며 진행된다. 각 phase에서
게임은 다음 \textit{안전구역}---팀이 머물 수 있는 영역---을 제시하고, 그
바깥의 \textit{블루존}은 밖에 있는 모두에게 점점 커지는 피해를
입힌다~\cite{pubgGameInfo}. 개발사 KRAFTON은 공식 developer API를 통해
match별 telemetry stream을 제공하며, 여기에는 일정 간격으로 기록되는
플레이어 위치와 시간 순서의 전투·구역·탈락 사건이
담긴다~\cite{pubgTelemetryEvents}. 본 연구의 분석 단위는 플레이어가 아니라
팀이다. 플레이어 수준 기록은 팀 수준의 trajectory, 노출, 팀 간 관계로
집계되며, 마지막 생존 팀원이 제거되는 시점을 팀 탈락 시점으로 정의한다.
이 stream은 본질적으로 시공간적이다. 팀은 이동 객체이고, 두 구역은 동적
공간 영역이며, 전투와 탈락은 특정 시점과 위치에 결합된 사건이다.

PUBG에 관한 대부분의 공개 benchmark와 예측 연구는 경기 후 집계로 final
placement를 예측한다~\cite{kagglePubgFinishPlacement,sen2021pubgFinalRank,lee2025pubgRank}.
이러한 문제 설정은 경기 단위의 최종 결과를 예측하는 데에는 적합하다. 다만
phase 수준의 생존 동역학을 다루도록 설계된 것이 아닐 뿐이다. 본 연구의
목적에 가장 가까운 것은 Sagong et al.~\cite{sagong2025hunterCamper}로,
PUBG 생존을 phase 단위로 분석하여 각 플레이어의 공격성과 활동성을
점수화하고 그로부터 얻은 전략 유형으로 next-phase 생존을 예측한다. 본
연구는 그들의 phase 수준 문제 설정을 이어받되, 분석 단위와 표현을 바꾼다.
분석 단위를 플레이어가 아니라 팀으로 두고, 소수의 전략 점수 대신 전체
telemetry를 phase에 정렬된 시공간 표현---팀 단위 feature와 동적 team
graph---으로 재구성한다. 이를 통해 다음 phase 생존 여부의 분류 자체를
넘어, 팀의 탈락 위험이 어떤 시공간 조건과 팀 간 관계 속에서 형성되는지를
분석한다.

구체적으로, 본 연구는 각 phase의 종료 시점까지 관측 가능한 정보로 그 팀이
다음 phase 동안 탈락할 확률을 예측한다. 이 양을 \textit{next-phase
elimination hazard}라고 부른다. 이는 생존을 보는 이산 시간 event-history
관점이며~\cite{allison1982discrete,singer1993discrete}, 하나의 match 수준
label을 팀과 phase마다 하나씩의 위험 label로 대체한다. 이 과제를 지탱하는
일은 feature engineering 이상을 요구한다. Feature, label, 관계가 모두 같은
시간·공간 축 위에서 phase 단위로, 미래 정보를 건드리지 않고 도출되어야
한다. 따라서 telemetry를 phase-정렬 시공간 데이터베이스(STDB)로
재구성하며, 그 객체들---이동하는 팀, 동적 구역, 공간 사건, phase
interval---이 그러한 도출을 체계적이고 재현 가능하게 만든다. 마지막으로,
팀의 위험은 자신의 상태만이 아니라 주변 팀들의 배치, 즉 누가 가까이
있는지, 누가 누구를 공격하는지, 누가 같은 축소 공간으로 밀려나는지에
의해서도 결정된다. 본 연구는 이 관계를 phase마다 하나의 동적 team graph로
표현하고, 관계 자체가 고유한 신호를 갖는지 검증한다.

본 연구는 다음 두 연구 질문을 중심으로 구성된다.
\begin{description}
  \item[\textbf{RQ1.}] 공식 PUBG telemetry를 phase-정렬 STDB로 재구성하여, 팀의 trajectory, 구역 노출, 전투 사건, 동적 팀 간 관계를 team-phase feature와 dynamic team graph로 도출할 수 있는가?
  \item[\textbf{RQ2.}] 이 표현 위에서, (i) team-phase feature만으로 next-phase elimination hazard를 얼마나 정확히 예측할 수 있고, (ii) dynamic team graph는 그 feature만 사용하는 모델을 넘어 추가 예측 신호를 제공하는가?
\end{description}

본 연구의 기여는 세 가지다. 첫째, 배틀로얄 생존 예측을 final-placement
추정에서 팀 수준·phase 수준의 elimination hazard로 재정식화한다. 둘째, 이
과제를 성립시키는 표현을 제공한다. 공식 developer API에서 수집한 21,824개
match로 구축한 phase-정렬 STDB와, 누수 없이 도출된 team-phase feature 및
dynamic team graph이다. 셋째, 예측 정보가 어디에 있는지 정량화한다. 선행
연구의 문제 설정을 반영한 strategy-score baseline을 전체 telemetry 기반
표현으로 대체하면, pooled aligned 평가에서 동일한 학습 알고리즘의 AUC가
0.649에서 0.765로 상승한다. 또한 dynamic team graph는 team-phase feature만 사용하는 모델
위에 추가 향상을 제공하며, 이 효과는 match 단위 paired test, edge
ablation, 무작위화 통제로 확인된다. 정확도 너머에서, phase 수준의 관점은
팀의 위험이 언제 어떤 공간 조건에서 상승하는지를 위치시키며, phase 경계마다 갱신되어
위험에 처한 팀을 식별하는 중계 분석, 경쟁 플레이를 위한 phase별 피드백, 그리고
live-service 게임 telemetry를 생존 과정으로 분석하는 후속 연구의 출발점이
될 수 있다.
